\section{Objetivos}

Os objetivos do trabalho estão organizados em um objetivo geral e objetivos específicos.

\subsection{Objetivo geral}
Desenvolver uma aplicação \textit{web} com o \textit{framework} Laravel para apoiar a gestão operacional de uma lanchonete, contemplando o cadastro e a autenticação de usuários, o gerenciamento de produtos, a exibição do catálogo para clientes e o processamento de pedidos, com persistência dos dados em MySQL.

\subsection{Objetivos específicos}

\begin{itemize}
    \item Levantar os requisitos funcionais e não funcionais do sistema de gestão de lanchonete.
    \item Modelar o banco de dados MySQL para armazenar informações de usuários, produtos e pedidos.
    \item Desenvolver um mecanismo de autenticação e controle de acesso para funcionários e clientes.
    \item Implementar o cadastro, edição, listagem e exclusão de produtos.
    \item Desenvolver a funcionalidade de exibição de produtos na área do cliente.
    \item Implementar o módulo de carrinho de compras e processamento de pedidos.
    \item Gerenciar pedidos e permitir a atualização de seus estados para pendente, em preparo ou finalizado.
    \item Integrar a aplicação ao banco de dados MySQL, garantindo a persistência das informações.
    \item Verificar os principais fluxos funcionais com testes automatizados de ponta a ponta e analisar a percepção de usabilidade dos participantes por meio de questionário pós-tarefa.
\end{itemize}
